% -------------------------------------------------------------------------
% WBS ID: RES-NUM-MSH
% Deliverable: Mesh Topology, Resolution and Geometric Discretisation
% Status: Draft complete
% Evidence status: Regional and local mesh requirements established at research level
% Review status: Not reviewed
% Last updated: 2026-07-18
% Dependencies: RES-NUM-FRM, RES-DAT-GEO, RES-GEO-REP, RES-MOD-SYS
% Downstream interfaces: RES-NUM-FLX, RES-NUM-TIM, RES-VER-CNV, Software mesh implementation
% -------------------------------------------------------------------------

\section{RES-NUM-MSH --- Mesh Topology, Resolution and Geometric Discretisation}
\label{sec:res-num-msh}

\subsection{Purpose and Scope}
\label{sec:res-num-msh-purpose}

% TODO: Define what this Deliverable covers and explicitly excludes.

\subsection{Research Questions}
\label{sec:res-num-msh-research-questions}

% TODO: State the questions that must be answered before the Deliverable
% can be accepted.

\subsection{Terminology and Definitions}
\label{sec:res-num-msh-definitions}

% TODO: Define the technical terminology, symbols, quantities and
% classifications used in this Deliverable.

\subsection{Literature Evidence}
\label{sec:res-num-msh-literature-evidence}

% TODO: Record evidence from peer-reviewed papers, authoritative datasets,
% standards, technical reports and primary documentation.
%
% Do not create a detached literature-summary list. Explain how each source
% supports, contradicts or limits a project decision.

\textcite{roenby2016computational} and \textcite{maric2013vofoam}
support geometric VOF advection on general or unstructured
finite-volume meshes. Source role: `3D-VOF`. Project conclusion: the
local 3D mesh may use polyhedral finite volumes and refine around the
air--water interface and barrier interaction region. Limitation: the
sources do not set the final project cell sizes, refinement ratios,
wall-normal spacing or adaptive-mesh triggers.

\subsection{Governing Relations and Quantitative Definitions}
\label{sec:res-num-msh-governing-relations}

% TODO: Record relevant equations, dimensionless groups, empirical models,
% extraction rules or quantitative definitions.
%
% Use align environments for displayed equations.
% Every displayed equation must have a unique \label{eq:...}.
% Do not place punctuation after displayed equations.

\subsection{Discrete Formulation}
\label{sec:res-num-msh-discrete-formulation}

The regional finite-volume mesh is specified as a set of straight-edged
polygonal control volumes. The practical production mesh is unstructured
triangular; quadrilateral meshes are retained for canonical tests and
simple offshore corridors. The cell update is shape-independent and uses
only the geometry in
\cref{eq:res-num-spec-regional-cell-geometry-set}.

The local finite-volume mesh is specified as a three-dimensional
polyhedral mesh with cell volumes, face area vectors, owner--neighbour
connectivity, boundary patches and wall-distance information. The local
mesh shall represent the fixed terrain and barrier surfaces while
allowing refinement near the free surface and local impact region.

\subsection{Conservation Properties}
\label{sec:res-num-msh-conservation-properties}

% TODO: Complete this RES-NUM Domain-specific subsection.

\subsection{Consistency and Formal Accuracy}
\label{sec:res-num-msh-consistency-formal-accuracy}

% TODO: Complete this RES-NUM Domain-specific subsection.

\subsection{Stability Requirements}
\label{sec:res-num-msh-stability-requirements}

% TODO: Complete this RES-NUM Domain-specific subsection.

\subsection{Mesh and Time-Step Requirements}
\label{sec:res-num-msh-mesh-time-step-requirements}

Mesh quality directly affects gradient reconstruction, limiter
activation, hydrostatic source balance, wet/dry stability and the CFL
timestep in \cref{eq:res-num-spec-regional-cell-timestep}. The mesh
shall therefore record cell area, face length, outward normal, face
centroid, cell centroid and neighbour connectivity for every regional
cell.

For the local three-dimensional URANS--VOF model, refinement shall be
considered in the following zones:

\begin{itemize}
  \item incident-wave inlet and relaxation region;
  \item expected free-surface path and interface-fragmentation region;
  \item run-up and overtopping paths;
  \item barrier leading edges, corners, crest and openings;
  \item near-wall terrain and barrier boundary layers;
  \item pressure-probe and force-integration surfaces;
  \item reflected and transmitted wave measurement locations;
  \item downstream wake and recirculation region;
  \item outlet and side-boundary absorption regions.
\end{itemize}

The local mesh-quality record shall include non-orthogonality, skewness,
aspect ratio, local cell volume, face area, wall distance, local
Courant number, interface Courant number and wall-function validity
indicators. These quantities are numerical-quality outputs rather than
case-study conclusions.

\subsection{Numerical Failure Modes}
\label{sec:res-num-msh-numerical-failure-modes}

% TODO: Complete this RES-NUM Domain-specific subsection.

\subsection{Verification Requirements}
\label{sec:res-num-msh-verification-requirements}

% TODO: Complete this RES-NUM Domain-specific subsection.

\subsection{Candidate Approaches}
\label{sec:res-num-msh-candidate-approaches}

% TODO: Define the credible candidate approaches considered.

\subsection{Critical Comparison}
\label{sec:res-num-msh-critical-comparison}

% TODO: Compare candidates using physically and project-relevant criteria.
% Do not select an approach solely on popularity or implementation ease.

\subsection{Assumptions and Applicability}
\label{sec:res-num-msh-assumptions}

% TODO: State assumptions and the conditions under which they are valid.

\subsection{Limitations and Failure Conditions}
\label{sec:res-num-msh-limitations}

% TODO: State where the evidence, model or selected approach ceases to be
% reliable or applicable.

\subsection{Project-Specific Conclusions}
\label{sec:res-num-msh-conclusions}

% TODO: Convert the evidence into conclusions specific to the Studentship
% project. Do not merely restate literature findings.

\subsection{Selected Approach and Rationale}
\label{sec:res-num-msh-selected-approach}

% TODO: Record the selected approach only after sufficient evidence exists.
% State the alternatives rejected and the reasons for rejection.

The selected regional mesh approach is:

\begin{itemize}
  \item general polygonal finite-volume equations;
  \item unstructured triangles as the primary case-study mesh;
  \item quadrilaterals for verification, simple corridors and structured
    comparisons;
  \item refinement around bathymetric gradients, shorelines, wet/dry
    fronts, source regions, candidate defences, gauge/observation
    locations and regional-to-local extraction boundaries.
\end{itemize}

The selected local mesh approach is:

\begin{itemize}
  \item body-fitted or otherwise geometry-resolving polyhedral
    finite-volume cells for the fixed terrain and fixed barrier;
  \item stable boundary patches matching the local boundary-role
    partition defined in `RES-MOD-IBC`;
  \item local refinement around the free surface, impact surfaces,
    overtopping region, wake and validation probes;
  \item wall-normal resolution sufficient for the selected high-Re
    wall-function treatment;
  \item mesh-convergence studies based on free-surface, pressure,
    force, moment and overtopping metrics.
\end{itemize}

\subsection{Unresolved Decisions}
\label{sec:res-num-msh-unresolved-decisions}

% TODO: Record unresolved questions, required evidence and decision owners.

The mesh topology is specified at research level. Final mesh sizes,
quality thresholds, refinement ratios and QGIS-derived domain extents
remain open until `RES-DAT-GEO`, `RES-GEO-REP` and validation planning
close the case-study geometry.

For the local model, final wall-normal spacing, roughness treatment,
static-versus-adaptive refinement policy, interface-refinement
thresholds and domain extents remain open until the case-study defence
geometry and validation data are locked.

\subsection{Requirements and Handoffs}
\label{sec:res-num-msh-requirements}

% TODO: State requirements passed to other Research Domains, Software,
% verification activities or later execution work.

`RES-NUM-ERR` shall record mesh-quality and refinement indicators for
local simulations. `RES-VER-CNV` shall define local mesh-convergence
tests for free-surface elevation, pressure, force, moment, overtopping
volume and energy metrics.

\subsection{Deliverable Completion Record}
\label{sec:res-num-msh-completion-record}

% TODO: Record whether the Deliverable is:
%
% - Not started
% - In progress
% - Draft complete
% - Reviewed
% - Accepted
%
% Acceptance requires:
%
% - the research questions to be answered;
% - claims to be supported by appropriate evidence;
% - assumptions and limitations to be explicit;
% - the project-specific conclusion to be stated;
% - downstream requirements to be traceable;
% - unresolved decisions to be closed or formally deferred.
